We study an exothermic pseudo-Dirac dark sector in which a vector mediator couples to quarks with strengths proportional to their masses. A reference xenon recoil yield and the Planck relic abundance determine the dark mass and coupling once the mediator mass and the top coupling are chosen. Scanning fourteen mediator masses between 20 and 200~GeV gives 470 admissible points, each generated separately and passed through the CMS all-hadronic $t\bar t$+MET selection. A 15-bin CLs test at $138\ \mathrm{fb}^{-1}$ nominally excludes 381 of them at 95\% CL and leaves 89; every point at or above 100~GeV is excluded. Between 50 and 90~GeV the excluded region runs from a lower boundary at $g_t\approx0.17$ to the top of the scanned range, at 30 and 40~GeV a second boundary closes it at $g_t\approx2.7$ and $3.0$, and at 20~GeV a single point is nominally excluded within the Monte Carlo uncertainty. Three levels of test support the map: the relic equation solved at both Planck-window edges for all 470 nodes, an approximate rescaling of the signal for the same nodes, and dedicated endpoint samples at four nodes, which show that the efficiency carries no systematic shift beyond its statistical scatter. One node crosses $\mu_{95}=1$ at the window edges, by $0.05\%$. The analysis is an expected-only recast, and its limitations are stated explicitly.
